\documentclass[10pt,a4paper]{amsart}
\usepackage[margin=19mm]{geometry}
\usepackage{amsmath,amssymb,mathtools}
\usepackage[scr=rsfs,cal=euler]{mathalfa}
\usepackage{xfrac}
\usepackage[unicode,hidelinks,bookmarks=false]{hyperref}
\newcommand{\ie}{i.e.\ }
\newcommand{\cf}{cf.\ }
\newcommand{\resp}{respectively\ }
\newcommand{\Subsection}{\subsection}
\providecommand{\nobreakdash}{\nobreak\hbox{-}}
\setlength{\emergencystretch}{2em}
\multlinegap0pt
\usepackage{bm}
\usepackage{bbm}
\usepackage{dsfont}
\usepackage{xspace}
\usepackage{cancel}
\usepackage{mathtools}
\usepackage{amsthm}
\makeatletter
\@ifundefined{cor}{\newtheorem{cor}{Corollary}}{}
\@ifundefined{thm}{\newtheorem{thm}{Theorem}}{}
\makeatother
\title[The Kaczmarz Algorithm in Hilbert $C^{*}$-modules]{The Kaczmarz Algorithm in Hilbert $C^{*}$-modules}
\author{Daniel Alpay, Chad Berner, Eric S. Weber}
\date{Source: arXiv:2508.13861v1 (CC BY 4.0)}
\begin{document}
\maketitle

\noindent\emph{Source and licence.} The Kaczmarz Algorithm in Hilbert $C^{*}$-modules. Version arXiv:2508.13861v1 (CC BY 4.0), distributed under the Creative Commons Attribution 4.0 International licence, as recorded in the arXiv licence metadata for this exact version and shown on the abstract page https://arxiv.org/abs/2508.13861v1. Full rights notice: licenses/arxiv-2508.13861-CC-BY-4.0.txt.\par\smallskip

\noindent\textit{Editorial scope.} Counted unit: the Cor whose source label is \texttt{equivalenttoauxcor} in the article (source0.tex lines 545--554, source label \texttt{equivalenttoauxcor}), delivered together with the article's own complete original proof of it (source0.tex lines 555--583, one \texttt{proof} environment of 29 lines). No mathematics is written, repaired or re-derived: statement and proof are the article's own text line for line. Required context retained uncounted: imageNCT (lines 517--519). Every definition, display and label of those lines is retained unchanged. Omitted from this package and not redistributed: 1--516, 520--544, 584--609. Nothing inside a retained excerpt is omitted or altered: every retained body is delivered exactly as the article's lines are written, so no omission receipt of any kind accompanies this package. Citations: the retained excerpts carry no citation command at all, so no bibliography entry of the article is transcribed and no reference is redistributed. Reference adaptation: none was needed and none was made; every reference inside the delivered unit resolves inside it, so no unresolved reference or citation is produced. Printed slips kept literal: no printed word, symbol, hypothesis or number of the counted statement or of its proof was altered. Layout and class: the article's own class is \texttt{amsart} with packages including amsmath, amssymb, graphicx, biblatex; the wrapper is the lane's pinned \texttt{amsart} editorial wrapper at 10pt on A4, which restates the theorem-like environments the retained text needs and adds the article's own macro definitions that the retained text needs (none), with extra wrapper packages (bm, bbm, dsfont, xspace, cancel, mathtools). The delivered numbering is the unit's own. Assets: no retained span contains a figure environment, a graphics-inclusion command, a PDF-page inclusion, a table or a TikZ picture; no image file is copied into this package, the package declares no asset dependency and no image file is redistributed with it. Licence: version arXiv:2508.13861v1 of this article is distributed under the Creative Commons Attribution 4.0 International licence, as recorded in the arXiv licence metadata for this exact version and shown on the abstract page https://arxiv.org/abs/2508.13861v1. A full rights notice is delivered at licenses/arxiv-2508.13861-CC-BY-4.0.txt. Characters: every retained excerpt is pure ASCII, so the delivered engine, XeLaTeX with its default Latin Modern text font, reports no missing character.
\begin{thm}\label{imageNCT}
$$ImV_{\mu_{\{x\}}}=[b^{x}(w)C(X, H^{2}(\mathbb{D}))]^{\perp}$$ where $\{b^{x}(w)\}_{x\in X}$ is the family of inner functions and zero functions corresponding to $\{\mu_{x}\}_{x\in X}$ via the Herglotz theorem.
\end{thm}

\begin{cor}\label{equivalenttoauxcor}
Let $\{f_{n}\}_{n=0}^{\infty}$ be a standard frame in a Hilbert $C(X)$-module $H$ where $X$ is a compact Hausdorff space. The following are equivalent:

\begin{enumerate}
    \item $$Im\theta_{f_{n}}=[b^{x}(w)C(X, H^{2}(\mathbb{D}))]^{\perp}$$ where for each $x\in X$, $b^{x}(w)$ is an inner function such that $$b^{x}(0)=0$$ or $$b^{x}(w)=0,$$ and the family of measures corresponding to these functions via the Herglotz theorem $\{\mu_{x}\}_{x\in X}$ are weakly continuous.
    \item There is a bounded, invertible, operator with an invertible adjoint that sends $\{f_{n}\}_{n=0}^{\infty}$ to the auxiliary sequence of effective sequence $\{e^{2\pi i ny}\}_{n=0}^{\infty}\subseteq WC(X, L^{2}(\mu_{x}))$.
\end{enumerate}


\end{cor}

\begin{proof}
$(1)\implies (2)$

By Theorem \ref{imageNCT}, for all $f\in H$, there is a $B(f)\in WC(X, L^{2}(\mu_{x}))$ such that
\begin{equation}\label{ipeq}
\langle f, f_{n}\rangle_{H} =\langle B(f), g_{n}\rangle_{WC(X, L^{2}(\mu_{x}))} 
\end{equation} for all $n$ where $\{g_{n}\}$ is the auxiliary sequence of $\{e^{2\pi i ny}\}\subseteq WC(X, L^{2}(\mu_{x}))$. Also, clearly the function $B: H\to WC(X, L^{2}(\mu_{x}))$ is invertible. Now note that since $\{g_{n}\}$ is a standard Parseval frame and $\{f_{n}\}$ is a stadnard frame that $B$ is $A$-linear and bounded.

Now define $B^{*}: WC(X, L^{2}(\mu_{x}))\to H$ where $B^{*}(g_{n})=f_{n}$ for all $n$ and extend $C(X)$-linearly. To see that $B^{*}$ is well defined,
suppose that
$$\sum_{n=0}^{k}c_{n}f_{n}=0$$ for some $\{c_{n}\}\subseteq C(X)$. By equation \ref{ipeq} since $B$ is surjective,
$$\sum_{n=0}^{k}c_{n}g_{n}=0.$$ Therefore, $B^{*}$ is well defined on $span_{C(X)}\{g_{n}\}$.
Furthermore for any $\{c_{n}\}\subseteq C(X)$ by equation \ref{ipeq},
$$\left \langle \sum_{n=0}^{k}c_{n}f_{n}, \sum_{n=0}^{k}c_{n}f_{n} \right\rangle_{H}=\left\langle B\sum_{n=0}^{k}c_{n}f_{n},\sum_{n=0}^{k}c_{n}g_{n} \right\rangle_{WC(X, L^{2}(\mu_{x}))}.$$
Thus, $B^{*}$ is also bounded on $span_{C(X)}\{g_{n}\}$ since $B$ is bounded, showing $B^{*}$ uniquely extends continuously to $WC(X, L^{2}(\mu_{x}))$.

Now suppose that $$B^{*}f=0=B^{*}\sum_{n=0}^{\infty}\langle f, g_{n}\rangle_{WC(X, L^{2}(\mu_{x}))}g_{n}=\sum_{n=0}^{\infty}\langle f, g_{n}\rangle_{WC(X, L^{2}(\mu_{x}))}f_{n}=$$
$$\sum_{n=0}^{\infty}\langle B^{-1}f, f_{n}\rangle_{H}f_{n}.$$
Since $\{f_{n}\}$ is a frame, we have $f=0$.

Now to see that $B^{*}$ is surjective, note that for any $f\in H$, there is a $\{c_{n}\}\in Im\theta_{f_{n}}$ such that
$$\sum_{n=0}^{\infty}c_{n}f_{n}=f=B^{*}\sum_{n=0}^{\infty}c_{n}g_{n}.$$ The above equality is sound because $\sum_{n=0}^{\infty}c_{n}g_{n}$ converges for any $\{c_{n}\}\in Im\theta_{g_{n}}$ by the Parseval frame condition.

Now it is clear that $B^{*}$ is the adjoint of $B$.

$(2)\implies (1)$

This follows easily from Theorem \ref{imageNCT}.
\end{proof}

\end{document}
